% Generated draft: requires mathematical and source review.
\documentclass[11pt]{article}
\usepackage[margin=25mm]{geometry}
\usepackage{fontspec}
\setmainfont{Latin Modern Roman}
\usepackage{amsmath,amssymb,mathtools,unicode-math}
\setmathfont{Latin Modern Math}
\usepackage{enumitem,needspace,xurl,hyperref}
\setlength{\parindent}{0pt}
\setlength{\parskip}{5pt}
\setlength{\emergencystretch}{3em}
\newcommand{\N}{\mathbb{N}}
\newcommand{\Z}{\mathbb{Z}}
\newcommand{\Q}{\mathbb{Q}}
\newcommand{\R}{\mathbb{R}}
\newcommand{\C}{\mathbb{C}}
\newcommand{\F}{\mathbb{F}}
\newcommand{\A}{\mathbb{A}}
\newcommand{\PP}{\mathbb{P}}
\newcommand{\E}{\mathbb{E}}
\newcommand{\eps}{\varepsilon}
\newcommand{\dd}{\,\mathrm{d}}
\begin{document}
\section*{UnsolvedMath ID: 1297}
\textbf{Source problem code:} catalogue entry
\subsection{Definitions and mathematical statement}
Let $F$ be a totally real number field, meaning a finite extension of $\Q$ all of whose embeddings into $\C$ have image in $\R$, and let $\mathcal O_F$ be its ring of integers. For a commutative ring $A$, let $E(A)$ be the stable elementary matrix group, generated by the matrices $I+ae_{ij}$, with $a\in A$ and $i\ne j$, inside the direct limit of the general linear groups. The stable Steinberg group $\operatorname{St}(A)$ has generators $x_{ij}(a)$ and Steinberg relations
\[
x_{ij}(a)x_{ij}(b)=x_{ij}(a+b),
\]
\[
[x_{ij}(a),x_{kl}(b)]=1\quad(i\ne l,\ j\ne k),\qquad
[x_{ij}(a),x_{jk}(b)]=x_{ik}(ab)\quad(i,j,k\text{ distinct}),
\]
where $[x,y]=xyx^{-1}y^{-1}$. Its map to $E(A)$ sends $x_{ij}(a)$ to $I+ae_{ij}$, and the stable group $K_2(A)$ is its kernel. This kernel is the appropriate $K$-theoretic object; an unspecified finite-rank Steinberg-group centre is not used in its place.

Let $\zeta_F(s)$ be the Dedekind zeta function, initially
\[
\zeta_F(s)=\sum_{0\ne\mathfrak a\subset\mathcal O_F}N(\mathfrak a)^{-s}\qquad(\operatorname{Re}(s)>1),
\]
where the sum ranges over nonzero integral ideals and $N(\mathfrak a)=|\mathcal O_F/\mathfrak a|$; its value at $s=-1$ is defined by analytic continuation. Define $w_2(F)$ as the greatest positive integer $m$ for which every $\sigma\in\operatorname{Gal}(\overline F/F)$ satisfies
\[
\chi_m(\sigma)^2\equiv1\pmod m,
\]
where $\chi_m$ is the cyclotomic character describing the action of $\sigma$ on the $m$th roots of unity. Equivalently, $w_2(F)$ is the order of the Galois invariants of the second Tate twist $\Q/\Z(2)$; the cyclotomic description specifies the normalization.

The intended standard Birch--Tate formula recovered by the supplied triage is
\[
\forall\text{ totally real number fields }F,\qquad
|K_2(\mathcal O_F)|=w_2(F)\,|\zeta_F(-1)|.
\]
The source's literal question omits the totally real hypothesis, the evaluation point, and this normalization. Its background mentions $s=-1$ but gives an unnormalized centre equality; the supplied triage warns that the unnormalized equality is false even for $\Q$ and that identifying a Steinberg-group centre requires stability hypotheses. The displayed formula is thus an explicitly qualified restoration, not the literal background assertion.

The registry labels the record open. Its literature triage dated 2026-08-17 classifies it as partially solved, attributes the odd-primary part for every totally real field to Wiles, and reports the full formula for totally real abelian extensions of $\Q$. It says that the $2$-primary part follows from the $2$-adic Iwasawa main conjecture and remains the outstanding part for arbitrary nonabelian totally real fields. In valuation terms, the reported odd-primary progress concerns every odd prime $\ell$ in
\[
v_\ell(|K_2(\mathcal O_F)|)=v_\ell(w_2(F)|\zeta_F(-1)|),
\]
where $v_\ell$ is extended to nonzero rational numbers in the usual way. These are attributed dated assessments, not a new verification of the literature.
\subsection{Short English statement}
For a totally real number field, the Birch--Tate formula predicts that the order of $K_2$ of its integers equals $w_2(F)$ times the absolute Dedekind zeta value at $-1$. The stored centre question omits essential hypotheses and normalization. The supplied review reports the remaining general issue as the nonabelian $2$-primary case.
\subsection{Sources}
\begin{itemize}
\item Charles A. Weibel, The K-book: An Introduction to Algebraic K-theory, Chapter VI, Section 8, Conjecture 8.6 and Theorem 8.7, supplied reference: \url{https://sites.math.rutgers.edu/~weibel/Kbook/Kbook.VI.pdf}. The record cites this for the normalized formula, Wiles's odd-primary result, and the abelian and remaining $2$-primary cases. The citation is retained without opening or generating a PDF.
\item Encyclopedia of Mathematics, ``Birch-Tate conjecture'': \url{https://encyclopediaofmath.org/wiki/Birch-Tate_conjecture}.
\item Local registry data/unsolvedmath/problems.json, permanent catalogue entry, source code catalogue entry, ``Birch--Tate Conjecture''; statement, background, and literature triage dated 2026-08-17. Dataset attribution: Ulam AI and the UnsolvedMath Contributors.
\end{itemize}
\end{document}
